\section{Formal proofs for exposure, history exclusion, and cycle walking}
\label{app:formal-proofs}

\subsection{Residual uniformity after observed credentials}

Let $D$ be a rank domain of size $N$ and let $\pi$ be uniformly sampled from
all permutations on $D$.  Condition on $q<N$ distinct consistent pairs
$(g_i,r_i)$, so $\pi(g_i)=r_i$ for $1\leq i\leq q$.  Let $g^*$ be an
unobserved input and let $r^*$ be any rank outside
$R=\{r_1,\ldots,r_q\}$.

\textbf{Theorem A.1 (observed-credential residual uniformity).}
Under the conditioning above,
\[
\Pr[\pi(g^*)=r^*\mid \pi(g_i)=r_i,1\leq i\leq q]
=\frac{1}{N-q}.
\]
Consequently $\Unrank_P(\pi(g^*))$ is uniform over the $N-q$ legal passwords
whose ranks have not been observed.

\textit{Proof.}
Exactly $(N-q)!$ permutations are consistent with the observed pairs.  Fixing
the additional pair $(g^*,r^*)$ leaves $(N-q-1)!$ consistent permutations.
Their ratio is $1/(N-q)$.  The rank/unrank bijection transfers this statement
to the accepted language. \hfill$\square$

This is standard random-permutation symmetry specialized to the credential
sequence, not a new combinatorial theorem.  For a real finite-domain PRP, let
$\mathcal A$ distinguish the unseen EPSCD rank from this conditional ideal
game after at most $q$ observed or chosen generation queries.  Replacing the
real permutation with a uniformly random one changes $\mathcal A$'s success
probability by at most the PRP advantage at the stated domain, query, and tweak
bounds.  Any preceding replacement of a Root-Key-derived credential key by an
independent key additionally incurs the relevant HKDF/PRF advantage.  The
statement does not turn the ideal equality into an unconditional claim about
the bounded FF1 prototype.

\subsection{Uniform exclusion}

Let $D=\{0,\ldots,N-1\}$ be the rank domain and let
$\Unrank:D\rightarrow\Lang_P$ be a bijection.  The authenticated history is
first rederived, restricted to $\Lang_P$, and deduplicated as a set of strings.
Its inverse rank set is
\[
B=\{\Rank_P(w):w\in E\cap\Lang_P\}\subset D,
\qquad e=|B|.
\]
Thus repeated history descriptors that rederive the same password contribute
one element to $B$, not multiple exclusions.

Fix distinct, unused generation inputs $g_1,g_2,\ldots$ in $D$, and sample
$\pi$ uniformly from the $N!$ permutations of $D$.  Write $Y_i=\pi(g_i)$ and
let
\[
T=\min\{i\geq 1:Y_i\notin B\}.
\]

\textbf{Theorem A.2 (uniform exclusion marginal).}
If $e<N$ and at least $e+1$ distinct generation inputs are available, then
$T\leq e+1$.  For every $z\in D\setminus B$,
\[
\Pr[Y_T=z]=\frac{1}{N-e},
\qquad
\mathbb{E}[T]=\frac{N+1}{N-e+1}.
\]
Consequently, $\Unrank(Y_T)$ is uniform on $\Lang_P\setminus E$.

\textit{Proof.}
For every $t\leq N$ and every tuple of distinct values
$(y_1,\ldots,y_t)$, a uniformly sampled permutation satisfies
\[
\Pr[(Y_1,\ldots,Y_t)=(y_1,\ldots,y_t)]=\frac{1}{(N)_t},
\]
where $(N)_t=N(N-1)\cdots(N-t+1)$.  Hence the images of the distinct scanned
inputs form an ordered sample without replacement.

No more than the $e$ distinct members of $B$ can occur before the first value
outside $B$, which proves $T\leq e+1$.  For a fixed $z\in D\setminus B$,
\[
\Pr[Y_T=z]
=\sum_{t=1}^{e+1}\frac{(e)_{t-1}}{(N)_t}.
\]
The right-hand side does not depend on $z$.  Since termination is certain and
the $N-e$ events $\{Y_T=z\}$ partition the sample space, each probability is
$1/(N-e)$.  Applying the rank/unrank bijection gives the password-space
marginal.

For the expectation, expose a uniformly random ordering of all $N$ ranks and
mark its $N-e$ non-excluded positions.  $T$ is the minimum marked position.
For $0\leq k\leq e$,
\[
\Pr[T>k]=\frac{(e)_k}{(N)_k}
=\frac{\binom{N-k}{N-e}}{\binom{N}{N-e}}.
\]
The tail-sum identity followed by the hockey-stick identity yields
\[
\mathbb{E}[T]
=\sum_{k=0}^{e}\Pr[T>k]
=\frac{1}{\binom{N}{N-e}}
  \sum_{j=N-e}^{N}\binom{j}{N-e}
=\frac{\binom{N+1}{N-e+1}}{\binom{N}{N-e}}
=\frac{N+1}{N-e+1}.
\tag*{$\square$}
\]

The theorem fixes $E$ independently of the newly sampled ideal permutation.
In the protocol, this means the history descriptors are authenticated and
their passwords are rederived before candidate scanning.  Undocumented remote
history and adversarial modification of local descriptors are outside this
probability statement.

\subsection{Cycle-walk termination and ideal tail}

Let $D=\{0,\ldots,N-1\}$ and embed it in
$U=\{0,\ldots,M-1\}$, where $M=2^{\lceil\log_2N\rceil}$.  Starting at
$x\in D$, apply a permutation $P$ on $U$ until its iterate returns to $D$.
Let $W$ count primitive permutation calls, including the accepting call.

\textbf{Lemma A.3 (deterministic termination).}
For every permutation $P$ and every $x\in D$, unbounded cycle walking stops
within $M-N+1$ primitive calls.

\textit{Proof.}
The orbit of $x$ is a finite cycle containing $x\in D$.  Before the next visit
to $D$, it cannot repeat an outside point and there are only $M-N$ such points.
The following call either reaches another member of $D$ or returns to $x$.
\hfill$\square$

\textbf{Lemma A.4 (ideal-permutation tail).}
If $P$ is a uniformly random permutation on $U$, then for every integer
$k\geq0$,
\[
\Pr[W>k]=\frac{(M-N)_k}{(M)_k},
\]
where the numerator is defined as zero for $k>M-N$.  If $N<M$, then
\[
\Pr[W>k]
\leq\left(\frac{M-N}{M}\right)^k<2^{-k}.
\]
If $N=M$, $W=1$ deterministically.

\textit{Proof.}
The event $W>k$ says that the first $k$ orbit images are distinct members of
$U\setminus D$.  Under a random permutation, these images are sampled without
replacement, giving
\[
\frac{M-N}{M}\frac{M-N-1}{M-1}\cdots
\frac{M-N-k+1}{M-k+1}=\frac{(M-N)_k}{(M)_k}.
\]
Each factor is at most $(M-N)/M$.  For a non-power-of-two $N$, the definition
of $M$ gives $M/2<N<M$, hence $(M-N)/M<1/2$.  The power-of-two case has no
outside point and accepts after one call. \hfill$\square$

For a cap of $c$ primitive calls, the ideal-model probability of a cap failure
is $\Pr[W>c]$, not $\Pr[W\geq c]$.  The concrete FF1 implementation is only
computationally related to the ideal permutation experiment.  The bound above
is therefore an ideal-model tail calculation; it is not an unconditional
availability guarantee for FF1 or its implementation.
